\begin{lemma}
\label{editorial-source-lemma-subalgebra-geometrically-connected}
Let $k$ be a field. Let $S$ be a $k$-algebra.
\begin{enumerate}
\item If $S$ is geometrically connected over $k$ so is every
$k$-subalgebra.
\item If all finitely generated $k$-subalgebras of $S$ are
geometrically connected, then $S$ is geometrically connected.
\item A directed colimit of geometrically connected $k$-algebras
is geometrically connected.
\end{enumerate}
\end{lemma}

\begin{proof}
Use throughout the nonzero-ring
criterion of Lemma
\ref{algebra-lemma-characterize-spec-connected}.
For (1), the original inclusion
$S'\subseteq S$ gives an injection
$S'\otimes_k L\to S\otimes_k L$
for every field $L/k$ by extension
of a vector-space basis. Its
target is nonzero, and the
unital source is therefore
nonzero as well. Any nontrivial
idempotent of the source would
remain nontrivial in the target,
which is impossible. This proves (1).

For (3), retain the original
nonempty directed system $(S_i)$
and its unital transition maps;
no injectivity is required.
For every $L/k$, the maps
$$
(\mathop{\rm colim}_i S_i)\otimes_k L
\longrightarrow\mathop{\rm colim}_i(S_i\otimes_k L),
\qquad[i,s]\otimes a\longmapsto[i,s\otimes a]
$$
and $[i,\sum_j s_j\otimes a_j]\mapsto\sum_j[i,s_j]\otimes a_j$
are well-defined inverses:
all representatives and balancing
equations involve finitely many
terms and can be compared at
one common later stage.
Each original $S_i\otimes_k L$
is nonzero. The colimit is
nonzero, since $1=0$ there
would hold at a later stage.
If $e$ is an idempotent in
the colimit, represent it at
a stage $i$. Its equation
$e^2-e=0$ holds at some later
stage $j$, making that later
representative idempotent.
It is $0$ or $1$ there by
connectedness, and hence is
$0$ or $1$ in the colimit.
This proves connectedness after
each $L/k$, and gives (3).

For (2), the finite type
$k$-subalgebras form a directed
system under adjoining the union
of two finite lists of generators,
and their union is the original
$S$. Apply (3). The hypotheses
exclude $S=0$, since its
subalgebra generated by the
image of $k$ would already
be zero and would fail
geometric connectedness.

Unlike geometric irreducibility,
geometric connectedness need
not survive a nonzero localization.
For the exact example
$$
S=k[x,y]/(xy),\qquad s=x+y,
$$
each field extension gives
two affine lines meeting at
their origin. More explicitly,
$V(x)$ and $V(y)$ cover its
spectrum, are irreducible
because their quotients are
$L[y]$ and $L[x]$, and meet
in the nonempty $\Spec(L)$.
Their union is therefore
connected. But in the original
$S_s$ the element
$$
e=\frac{x}{x+y},\qquad
e^2-e=\frac{-xy}{(x+y)^2}=0
$$
is a nontrivial idempotent.
The map
$$
S_s\longrightarrow k[x,x^{-1}]\times k[y,y^{-1}],
\qquad\frac{h(x,y)}{(x+y)^n}
\longmapsto\left(\frac{h(x,0)}{x^n},\frac{h(0,y)}{y^n}\right)
$$
sends it to $(1,0)$, so it
is neither $0$ nor $1$.
This map is in fact an
isomorphism with the following
explicit inverse. In $eS_s$,
whose unit is $e$, send $x$
to the original $x=ex$ and
$x^{-1}$ to $e/s$; their
product is $e$ because
$x(e/s)=x^2/s^2=e$.
In $(1-e)S_s$, send $y$ to
$y=(1-e)y$ and $y^{-1}$ to
$(1-e)/s$; their product is
$1-e$. Add the images of
the two Laurent polynomials.
Orthogonality of $e,1-e$
makes this a ring map.
The composites fix both
component units and Laurent
generators. On the original
side they fix $x,y$ and
$s^{-1}$, since
$e/s+(1-e)/s=s^{-1}$.
Thus both composites are
identities, proving the full
product description and the
failure of localization stability.
\end{proof}